\section{Dataflow and Precision Mechanisms}

\subsection{Coherent FP16}

A naive half-cache path stores RBF values in FP16 while constructing the local
$K[W,W]$ in FP32. The optimizer then consumes two inconsistent versions of
the same kernel. On EMBER-65K, this design requires 192 outer rounds, compared
with 152 for the direct FP32-cache control.

The coherent-FP16 path instead enforces the following execution invariant:
all decision-relevant local and global kernel values derive from the same
quantized inputs and distance rules. The implementation uses half precision
for training features or persistent RBF values, Tensor Cores for
high-dimensional dot-product tiles, and FP32 for the RBF transform,
accumulation, dual variables, gradients, and SMO control. It also avoids a
large FP32 staging panel.

With coherent values, outer rounds fall to 146. The mechanism is therefore not
simply that FP16 arithmetic is faster. A stateful optimizer must see one
coherent approximate operator at every point where it makes decisions.

\subsection{Shape-aware low-dimensional reuse}

The general high-dimensional TMA path pads SUSY from $d=18$ to $d=64$. The
compact path stores the original 18 features, uses aligned \code{half2} loads,
keeps two adjacent training rows per thread, stages 32 working-set rows at a
time, and accumulates RBF and gradient updates in FP32.

Relative to the strongest reuse-aware FP32 baseline, compact coherent FP16
improves SUSY-65K by 1.0995$\times$ and SUSY-1M by 1.0958$\times$. The generic
padded path is approximately neutral or slower. Low- and high-dimensional
workloads therefore require different kernels: Tensor Core dot products are
valuable in high dimensions, while eliminating padding and repeated row loads
matters more in low dimensions.

\subsection{Decision-path-aware block-solver micro-optimization}

The local solver repeatedly scans candidates, ranks lower variables, computes
an accepted step, and identifies threads that own selected variables. The
accepted micro-optimization uses approximate division only for broad
candidate ranking, retains exact FP32 arithmetic for the committed step $q$,
and replaces an all-thread owner search with direct index mapping.

Across eight captured snapshots, the candidate reproduces the captured
trajectory. In the formal end-to-end run, however, iteration counts differ
slightly (46,854 versus 46,861) and one of 65,536 training predictions differs.
The defensible description is therefore \emph{snapshot-equivalent and
quality-equivalent under the coherent-FP16 contract}, not globally
bitwise-exact.

\subsection{Lower precision beyond FP16}

Lower bit widths do not automatically extend the FP16 result.

\begin{itemize}
  \item An MXFP8-compute path with an FP16 cache retains much of the measured
        quality but does not establish a uniformly stronger mechanism.
  \item A linear U8 RBF cache doubles cache-column capacity and reduces one
        full-EMBER fit from \SI{4.865}{s} to \SI{3.734}{s}
        (1.303$\times$), but test accuracy falls from 0.92648 to 0.75329.
  \item An NVFP4 reconstructed-input gate reaches the internal tolerance and
        changes EMBER-65K training accuracy by $-3.05\times10^{-5}$. Its
        fit-time ratio is approximately 1.007$\times$ and preprocessing is
        excluded; the SVM kernel is not executed natively in FP4.
\end{itemize}

Feature quantization, RBF exponentiation, cache encoding, local
materialization, and solver state must form one coherent data path. Peak FP8
or FP4 GEMM throughput alone is not an end-to-end argument.
